% laborch package manual. Generated from docs/laborch_manual.md by build_manual.py
% (pandoc + LaTeX math and table-spacing fixes). Build: xelatex laborch_manual.tex (twice).
% Options for packages loaded elsewhere
\PassOptionsToPackage{unicode}{hyperref}
\PassOptionsToPackage{hyphens}{url}
\PassOptionsToPackage{dvipsnames,svgnames*,x11names*}{xcolor}
%
\documentclass[
  10pt,
]{article}
\usepackage{amssymb,amsmath}
\usepackage{ifxetex,ifluatex}
\ifnum 0\ifxetex 1\fi\ifluatex 1\fi=0 % if pdftex
  \usepackage[T1]{fontenc}
  \usepackage[utf8]{inputenc}
  \usepackage{textcomp} % provide euro and other symbols
\else % if luatex or xetex
  \usepackage{unicode-math}
  \defaultfontfeatures{Scale=MatchLowercase}
  \defaultfontfeatures[\rmfamily]{Ligatures=TeX,Scale=1}
  \setmainfont[]{DejaVu Serif}
  \setsansfont[]{DejaVu Sans}
  \setmonofont[Scale=0.85]{DejaVu Sans Mono}
\fi
% Use upquote if available, for straight quotes in verbatim environments
\IfFileExists{upquote.sty}{\usepackage{upquote}}{}
\IfFileExists{microtype.sty}{% use microtype if available
  \usepackage[]{microtype}
  \UseMicrotypeSet[protrusion]{basicmath} % disable protrusion for tt fonts
}{}
\makeatletter
\@ifundefined{KOMAClassName}{% if non-KOMA class
  \IfFileExists{parskip.sty}{%
    \usepackage{parskip}
  }{% else
    \setlength{\parindent}{0pt}
    \setlength{\parskip}{6pt plus 2pt minus 1pt}}
}{% if KOMA class
  \KOMAoptions{parskip=half}}
\makeatother
\usepackage{xcolor}
\IfFileExists{xurl.sty}{\usepackage{xurl}}{} % add URL line breaks if available
\IfFileExists{bookmark.sty}{\usepackage{bookmark}}{\usepackage{hyperref}}
\hypersetup{
  pdftitle={laborch: Package Manual},
  pdfauthor={LabOrch project},
  colorlinks=true,
  linkcolor=blue,
  filecolor=Maroon,
  citecolor=Blue,
  urlcolor=blue,
  pdfcreator={LaTeX via pandoc}}
\urlstyle{same} % disable monospaced font for URLs
\usepackage[margin=2.2cm]{geometry}
\usepackage{color}
\usepackage{fancyvrb}
\newcommand{\VerbBar}{|}
\newcommand{\VERB}{\Verb[commandchars=\\\{\}]}
\DefineVerbatimEnvironment{Highlighting}{Verbatim}{commandchars=\\\{\}}
% Add ',fontsize=\small' for more characters per line
\usepackage{framed}
\definecolor{shadecolor}{RGB}{248,248,248}
\newenvironment{Shaded}{\begin{snugshade}}{\end{snugshade}}
\newcommand{\AlertTok}[1]{\textcolor[rgb]{0.94,0.16,0.16}{#1}}
\newcommand{\AnnotationTok}[1]{\textcolor[rgb]{0.56,0.35,0.01}{\textbf{\textit{#1}}}}
\newcommand{\AttributeTok}[1]{\textcolor[rgb]{0.77,0.63,0.00}{#1}}
\newcommand{\BaseNTok}[1]{\textcolor[rgb]{0.00,0.00,0.81}{#1}}
\newcommand{\BuiltInTok}[1]{#1}
\newcommand{\CharTok}[1]{\textcolor[rgb]{0.31,0.60,0.02}{#1}}
\newcommand{\CommentTok}[1]{\textcolor[rgb]{0.56,0.35,0.01}{\textit{#1}}}
\newcommand{\CommentVarTok}[1]{\textcolor[rgb]{0.56,0.35,0.01}{\textbf{\textit{#1}}}}
\newcommand{\ConstantTok}[1]{\textcolor[rgb]{0.00,0.00,0.00}{#1}}
\newcommand{\ControlFlowTok}[1]{\textcolor[rgb]{0.13,0.29,0.53}{\textbf{#1}}}
\newcommand{\DataTypeTok}[1]{\textcolor[rgb]{0.13,0.29,0.53}{#1}}
\newcommand{\DecValTok}[1]{\textcolor[rgb]{0.00,0.00,0.81}{#1}}
\newcommand{\DocumentationTok}[1]{\textcolor[rgb]{0.56,0.35,0.01}{\textbf{\textit{#1}}}}
\newcommand{\ErrorTok}[1]{\textcolor[rgb]{0.64,0.00,0.00}{\textbf{#1}}}
\newcommand{\ExtensionTok}[1]{#1}
\newcommand{\FloatTok}[1]{\textcolor[rgb]{0.00,0.00,0.81}{#1}}
\newcommand{\FunctionTok}[1]{\textcolor[rgb]{0.00,0.00,0.00}{#1}}
\newcommand{\ImportTok}[1]{#1}
\newcommand{\InformationTok}[1]{\textcolor[rgb]{0.56,0.35,0.01}{\textbf{\textit{#1}}}}
\newcommand{\KeywordTok}[1]{\textcolor[rgb]{0.13,0.29,0.53}{\textbf{#1}}}
\newcommand{\NormalTok}[1]{#1}
\newcommand{\OperatorTok}[1]{\textcolor[rgb]{0.81,0.36,0.00}{\textbf{#1}}}
\newcommand{\OtherTok}[1]{\textcolor[rgb]{0.56,0.35,0.01}{#1}}
\newcommand{\PreprocessorTok}[1]{\textcolor[rgb]{0.56,0.35,0.01}{\textit{#1}}}
\newcommand{\RegionMarkerTok}[1]{#1}
\newcommand{\SpecialCharTok}[1]{\textcolor[rgb]{0.00,0.00,0.00}{#1}}
\newcommand{\SpecialStringTok}[1]{\textcolor[rgb]{0.31,0.60,0.02}{#1}}
\newcommand{\StringTok}[1]{\textcolor[rgb]{0.31,0.60,0.02}{#1}}
\newcommand{\VariableTok}[1]{\textcolor[rgb]{0.00,0.00,0.00}{#1}}
\newcommand{\VerbatimStringTok}[1]{\textcolor[rgb]{0.31,0.60,0.02}{#1}}
\newcommand{\WarningTok}[1]{\textcolor[rgb]{0.56,0.35,0.01}{\textbf{\textit{#1}}}}
\usepackage{longtable,booktabs}
% Correct order of tables after \paragraph or \subparagraph
\usepackage{etoolbox}
\makeatletter
\patchcmd\longtable{\par}{\if@noskipsec\mbox{}\fi\par}{}{}
\makeatother
% Allow footnotes in longtable head/foot
\IfFileExists{footnotehyper.sty}{\usepackage{footnotehyper}}{\usepackage{footnote}}
\makesavenoteenv{longtable}
\setlength{\emergencystretch}{3em} % prevent overfull lines
\providecommand{\tightlist}{%
  \setlength{\itemsep}{0pt}\setlength{\parskip}{0pt}}
\setcounter{secnumdepth}{-\maxdimen} % remove section numbering
\usepackage{fvextra}
\DefineVerbatimEnvironment{Highlighting}{Verbatim}{breaklines,breakanywhere,commandchars=\\\{\}}
\usepackage{etoolbox}
\AtBeginEnvironment{longtable}{\small}
\usepackage{titling}
\renewcommand{\arraystretch}{1.1}
\setlength{\emergencystretch}{3em}

\title{laborch: Package Manual}
\usepackage{etoolbox}
\makeatletter
\providecommand{\subtitle}[1]{% add subtitle to \maketitle
  \apptocmd{\@title}{\par {\large #1 \par}}{}{}
}
\makeatother
\subtitle{Version 0.2.0 (Phase 1)}
\author{LabOrch project}
\date{November 2026}

\begin{document}
\maketitle

{
\hypersetup{linkcolor=}
\setcounter{tocdepth}{2}
\tableofcontents
}
\texttt{laborch} is the Python core of LabOrch, a
computer-assisted orchestration tool. Given a
\textbf{Target} sound and an \textbf{Ensemble} of
instruments, it finds the combination of recorded notes
whose summed spectrum best matches the Target. The package
gathers the logic of the 2019--2020 notebooks (1.5, 1.7,
1.11, 1.12 and the \texttt{tool*.py} widget scripts) into
tested, importable modules (Phase 0), and turns the Bank
into a dataset (Phase 1): \textbf{Bank v2} holds power
spectra averaged over each note's automatically detected
sustain, keeps every recording's level, and comes with a
metadata table. See \texttt{docs/bank\_v2.md} for how it was
made and what it covers.

\subsection{1. Concepts and vocabulary}

\begin{longtable}[]{@{}lll@{}}
\toprule
\begin{minipage}[b]{0.30\columnwidth}\raggedright
Term\strut
\end{minipage} &
\begin{minipage}[b]{0.30\columnwidth}\raggedright
Meaning\strut
\end{minipage} &
\begin{minipage}[b]{0.30\columnwidth}\raggedright
In code\strut
\end{minipage}\tabularnewline
\midrule
\endhead
\begin{minipage}[t]{0.30\columnwidth}\raggedright
\textbf{Bank}\strut
\end{minipage} &
\begin{minipage}[t]{0.30\columnwidth}\raggedright
Every recorded note of every instrument (1,515 sounds, 33
instruments), each stored as an averaged spectrum\strut
\end{minipage} &
\begin{minipage}[t]{0.30\columnwidth}\raggedright
\texttt{Bank}\strut
\end{minipage}\tabularnewline
\begin{minipage}[t]{0.30\columnwidth}\raggedright
\textbf{Ensemble}\strut
\end{minipage} &
\begin{minipage}[t]{0.30\columnwidth}\raggedright
The instruments the composer writes for, e.g.~a string
quartet\strut
\end{minipage} &
\begin{minipage}[t]{0.30\columnwidth}\raggedright
a list of instrument codes,
\texttt{{[}"vln",\ "vla",\ "vlc"{]}}\strut
\end{minipage}\tabularnewline
\begin{minipage}[t]{0.30\columnwidth}\raggedright
\textbf{Palette}\strut
\end{minipage} &
\begin{minipage}[t]{0.30\columnwidth}\raggedright
Every Bank sound the Ensemble can play (notes, modes,
dynamics), as a matrix of spectra\strut
\end{minipage} &
\begin{minipage}[t]{0.30\columnwidth}\raggedright
\texttt{Palette} (from \texttt{Bank.select})\strut
\end{minipage}\tabularnewline
\begin{minipage}[t]{0.30\columnwidth}\raggedright
\textbf{Target}\strut
\end{minipage} &
\begin{minipage}[t]{0.30\columnwidth}\raggedright
The sound to imitate: a Bank sound or an uploaded
recording\strut
\end{minipage} &
\begin{minipage}[t]{0.30\columnwidth}\raggedright
a 2,048-element NumPy array\strut
\end{minipage}\tabularnewline
\begin{minipage}[t]{0.30\columnwidth}\raggedright
\textbf{Recommendation}\strut
\end{minipage} &
\begin{minipage}[t]{0.30\columnwidth}\raggedright
The weighted set of Palette sounds that best reproduces the
Target\strut
\end{minipage} &
\begin{minipage}[t]{0.30\columnwidth}\raggedright
\texttt{Fit}\strut
\end{minipage}\tabularnewline
\bottomrule
\end{longtable}

The model (Rose \& Hetrick 2009, ``RH2009'') treats the
Target spectrum $T$ as a weighted sum of Palette
spectra $X_j$:

\[
  T(f_i) \;\approx\; \sum_j X_{ij}\,\alpha_j, \qquad \alpha_j \ge 0
\]

The weights $\alpha$ are found by non-negative least squares
(NNLS), optionally weighted per frequency bin by $1/\sigma_i$ as in
\emph{Numerical Recipes} §15.4.

In Bank v2 the spectra are \textbf{power} spectra at each
recording's level. Powers of uncorrelated sounds add, so the
model is exact for a mix, and a weight αⱼ is a power gain on
recording j: αⱼ = 1 means ``as recorded'', αⱼ = 0.5 means 3
dB quieter. (In Bank v1 the spectra were unit-sum magnitudes
and a weight was only a share of the fit.)

\subsection{2. Installation}

Requirements: Python $\ge$ 3.9, NumPy, SciPy. Audio features
(reading recordings, splitting, mixing) also need
\textbf{pydub} and \textbf{ffmpeg} on the \texttt{PATH}.

From the \texttt{LabOrch} folder:

\begin{Shaded}
\begin{Highlighting}[]
\ExtensionTok{pip}\NormalTok{ install {-}e }\StringTok{".[dev]"}\NormalTok{     \# package + pydub + pytest, editable}
\CommentTok{\# or, without installing:}
\CommentTok{\#   import sys; sys.path.insert(0, "path/to/LabOrch")}
\end{Highlighting}
\end{Shaded}

On Windows, ffmpeg can be installed with
\texttt{winget\ install\ ffmpeg}.

The package finds the data relative to its own location (the
repository root). To use data elsewhere, set the environment
variable \texttt{LABORCH\_ROOT} before importing (see
\protect\hyperlink{51-laborchconfig}{5.1}).

\subsection{3. Quick start}

\begin{Shaded}
\begin{Highlighting}[]
\ImportTok{from}\NormalTok{ laborch }\ImportTok{import}\NormalTok{ Bank, nnls\_fit}

\NormalTok{bank }\OperatorTok{=}\NormalTok{ Bank.load()                     }\CommentTok{\# Bank v2: power spectra + metadata}
\NormalTok{target\_name }\OperatorTok{=} \StringTok{"trb{-}f2{-}ord{-}mf"}          \CommentTok{\# trombone F2, mf}
\NormalTok{target }\OperatorTok{=}\NormalTok{ bank.spectrum(target\_name)}

\NormalTok{palette }\OperatorTok{=}\NormalTok{ bank.select(instruments}\OperatorTok{=}\NormalTok{[}\StringTok{"clb"}\NormalTok{, }\StringTok{"bcl"}\NormalTok{],   }\CommentTok{\# Ensemble: clarinets}
\NormalTok{                      exclude}\OperatorTok{=}\NormalTok{[target\_name])}
\NormalTok{fit }\OperatorTok{=}\NormalTok{ nnls\_fit(palette, target)}

\BuiltInTok{print}\NormalTok{(fit.summary(}\DecValTok{4}\NormalTok{))}
\CommentTok{\# nnls: 62 sounds, relative error 0.644}
\CommentTok{\#     0.2184  clb{-}c5{-}ord{-}mf}
\CommentTok{\#     0.2065  clb{-}f4{-}ord{-}mf}
\CommentTok{\#     0.1714  clb{-}g4{-}ord{-}mf}
\CommentTok{\#     0.1095  clb{-}b4{-}ord{-}mf}

\NormalTok{best3 }\OperatorTok{=}\NormalTok{ fit.top(}\DecValTok{3}\NormalTok{, refit}\OperatorTok{=}\VariableTok{True}\NormalTok{)         }\CommentTok{\# keep 3 sounds, re{-}solve on them}
\NormalTok{best3.ranked()                         }\CommentTok{\# [(\textquotesingle{}clb{-}f4{-}ord{-}mf\textquotesingle{}, 0.346), ...]}

\NormalTok{bank.meta[target\_name][}\StringTok{"level\_db"}\NormalTok{]     }\CommentTok{\# {-}26.3 dBFS; the full metadata row}

\ImportTok{from}\NormalTok{ laborch.audio }\ImportTok{import}\NormalTok{ render\_fit   }\CommentTok{\# needs pydub + ffmpeg}
\NormalTok{render\_fit(bank, best3, n}\OperatorTok{=}\DecValTok{3}\NormalTok{, path}\OperatorTok{=}\StringTok{"trb\_from\_clarinets.wav"}\NormalTok{)}
\end{Highlighting}
\end{Shaded}

With every bin weighted equally ($\sigma = 1$), power spectra let
the strongest partials dominate, so the clarinets imitate
the trombone's loud upper partials rather than its
fundamental. Weighting by the Target's own spectrum
(\texttt{sigma=np.maximum(target,\ target.max()*1e-4)},
relative error per bin) brings \texttt{bcl-f2} and
\texttt{bcl-f3} back into the top four. Choosing $\sigma$ is Phase
2. With \texttt{Bank.load(version=1)} the example gives the
Phase 0 answer (\texttt{bcl-f2-ord-mf} first).

The notebook
\texttt{Notebooks/2.0\_laborch\_quickstart.ipynb} walks
through the same example with plots and playback.

\subsection{4. Package layout}

\begin{verbatim}
LabOrch/
├── laborch/
│   ├── __init__.py   public API (re-exports)
│   ├── config.py     paths and FFT constants
│   ├── bank.py       sound names, Bank (v2 and v1), Palette
│   ├── metadata.py   the metadata table: families, envelopes, ranges
│   ├── build.py      rebuild Bank v2 from the recordings
│   ├── spectra.py    audio loading, sustain detection, power and v1 spectra
│   ├── solve.py      DOT, pseudo-inverse and NNLS decompositions; Fit
│   ├── audio.py      gains and mixing for playback
│   └── synth.py      synthetic tones, spectrum-to-audio resynthesis
├── tests/            pytest suite (61 tests)
├── docs/             this manual, bank_v2.md report, make_bank_report.py
└── pyproject.toml
\end{verbatim}

Names importable directly from \texttt{laborch}:

\texttt{Bank}, \texttt{Palette}, \texttt{SoundName},
\texttt{parse\_name}, \texttt{INSTRUMENT\_NAMES},
\texttt{Fit}, \texttt{dot}, \texttt{nnls\_fit},
\texttt{pinv\_fit}, \texttt{chi2}, \texttt{make\_fft},
\texttt{power\_spectrum}, \texttt{find\_sustain},
\texttt{spectral\_peaks}, \texttt{FAMILIES}, \texttt{tone},
\texttt{spectrum\_to\_audio}.

Audio helpers are imported from their module
(\texttt{laborch.audio}, \texttt{laborch.spectra}) so that
the core works without pydub.

\begin{center}\rule{0.5\linewidth}{0.5pt}\end{center}

\subsection{5. Module reference}

\subsubsection{\texorpdfstring{5.1
\texttt{laborch.config}}{5.1 laborch.config}}

Constants only; no functions. Change paths here, not in
calling code.

\begin{longtable}[]{@{}lll@{}}
\toprule
\begin{minipage}[b]{0.30\columnwidth}\raggedright
Name\strut
\end{minipage} &
\begin{minipage}[b]{0.30\columnwidth}\raggedright
Value\strut
\end{minipage} &
\begin{minipage}[b]{0.30\columnwidth}\raggedright
Meaning\strut
\end{minipage}\tabularnewline
\midrule
\endhead
\begin{minipage}[t]{0.30\columnwidth}\raggedright
\texttt{ROOT}\strut
\end{minipage} &
\begin{minipage}[t]{0.30\columnwidth}\raggedright
repository folder, or \texttt{\$LABORCH\_ROOT}\strut
\end{minipage} &
\begin{minipage}[t]{0.30\columnwidth}\raggedright
base for every path below\strut
\end{minipage}\tabularnewline
\begin{minipage}[t]{0.30\columnwidth}\raggedright
\texttt{BANK\_DIR}\strut
\end{minipage} &
\begin{minipage}[t]{0.30\columnwidth}\raggedright
\texttt{ROOT/Soundbank\_McGill}\strut
\end{minipage} &
\begin{minipage}[t]{0.30\columnwidth}\raggedright
McGill Bank folder\strut
\end{minipage}\tabularnewline
\begin{minipage}[t]{0.30\columnwidth}\raggedright
\texttt{FFT\_DIR}\strut
\end{minipage} &
\begin{minipage}[t]{0.30\columnwidth}\raggedright
\texttt{BANK\_DIR/FFTBank}\strut
\end{minipage} &
\begin{minipage}[t]{0.30\columnwidth}\raggedright
current spectra, one \texttt{.txt} per sound\strut
\end{minipage}\tabularnewline
\begin{minipage}[t]{0.30\columnwidth}\raggedright
\texttt{SOUND\_DIRS}\strut
\end{minipage} &
\begin{minipage}[t]{0.30\columnwidth}\raggedright
\texttt{Soundbank\_WB}, \texttt{Soundbank\_Str},
\texttt{Soundbank\_PianoPercSax}\strut
\end{minipage} &
\begin{minipage}[t]{0.30\columnwidth}\raggedright
the \texttt{.aif} recordings\strut
\end{minipage}\tabularnewline
\begin{minipage}[t]{0.30\columnwidth}\raggedright
\texttt{OLD\_FFT\_DIR}\strut
\end{minipage} &
\begin{minipage}[t]{0.30\columnwidth}\raggedright
\texttt{ROOT/OldFFTBank}\strut
\end{minipage} &
\begin{minipage}[t]{0.30\columnwidth}\raggedright
2009 \texttt{.nbb} spectra (reference)\strut
\end{minipage}\tabularnewline
\begin{minipage}[t]{0.30\columnwidth}\raggedright
\texttt{MISC\_SOUNDS\_DIR}\strut
\end{minipage} &
\begin{minipage}[t]{0.30\columnwidth}\raggedright
\texttt{ROOT/MiscSounds}\strut
\end{minipage} &
\begin{minipage}[t]{0.30\columnwidth}\raggedright
extra Targets (trumpet mutes, multiphonics)\strut
\end{minipage}\tabularnewline
\begin{minipage}[t]{0.30\columnwidth}\raggedright
\texttt{CACHE\_FILE}\strut
\end{minipage} &
\begin{minipage}[t]{0.30\columnwidth}\raggedright
\texttt{ROOT/.cache/fftbank.npz}\strut
\end{minipage} &
\begin{minipage}[t]{0.30\columnwidth}\raggedright
binary cache of Bank v1\strut
\end{minipage}\tabularnewline
\begin{minipage}[t]{0.30\columnwidth}\raggedright
\texttt{BANK\_VERSION}\strut
\end{minipage} &
\begin{minipage}[t]{0.30\columnwidth}\raggedright
2\strut
\end{minipage} &
\begin{minipage}[t]{0.30\columnwidth}\raggedright
Bank loaded by default\strut
\end{minipage}\tabularnewline
\begin{minipage}[t]{0.30\columnwidth}\raggedright
\texttt{BANK\_FILE}\strut
\end{minipage} &
\begin{minipage}[t]{0.30\columnwidth}\raggedright
\texttt{BANK\_DIR/Bank\_v2/bank\_v2.npz}\strut
\end{minipage} &
\begin{minipage}[t]{0.30\columnwidth}\raggedright
Bank v2 names and power spectra (float32)\strut
\end{minipage}\tabularnewline
\begin{minipage}[t]{0.30\columnwidth}\raggedright
\texttt{BANK\_META}\strut
\end{minipage} &
\begin{minipage}[t]{0.30\columnwidth}\raggedright
\texttt{BANK\_DIR/Bank\_v2/bank\_v2.csv}\strut
\end{minipage} &
\begin{minipage}[t]{0.30\columnwidth}\raggedright
Bank v2 metadata table\strut
\end{minipage}\tabularnewline
\begin{minipage}[t]{0.30\columnwidth}\raggedright
\texttt{BUILD\_CACHE}\strut
\end{minipage} &
\begin{minipage}[t]{0.30\columnwidth}\raggedright
\texttt{ROOT/.cache/build\_v2.npz}\strut
\end{minipage} &
\begin{minipage}[t]{0.30\columnwidth}\raggedright
resumable state of \texttt{laborch.build}\strut
\end{minipage}\tabularnewline
\begin{minipage}[t]{0.30\columnwidth}\raggedright
\texttt{SAMPLE\_RATE}\strut
\end{minipage} &
\begin{minipage}[t]{0.30\columnwidth}\raggedright
44100\strut
\end{minipage} &
\begin{minipage}[t]{0.30\columnwidth}\raggedright
Hz\strut
\end{minipage}\tabularnewline
\begin{minipage}[t]{0.30\columnwidth}\raggedright
\texttt{NFFT}\strut
\end{minipage} &
\begin{minipage}[t]{0.30\columnwidth}\raggedright
4096\strut
\end{minipage} &
\begin{minipage}[t]{0.30\columnwidth}\raggedright
Hamming window length (samples)\strut
\end{minipage}\tabularnewline
\begin{minipage}[t]{0.30\columnwidth}\raggedright
\texttt{HOP}\strut
\end{minipage} &
\begin{minipage}[t]{0.30\columnwidth}\raggedright
512\strut
\end{minipage} &
\begin{minipage}[t]{0.30\columnwidth}\raggedright
step between windows (samples)\strut
\end{minipage}\tabularnewline
\begin{minipage}[t]{0.30\columnwidth}\raggedright
\texttt{NBINS}\strut
\end{minipage} &
\begin{minipage}[t]{0.30\columnwidth}\raggedright
2048\strut
\end{minipage} &
\begin{minipage}[t]{0.30\columnwidth}\raggedright
bins kept per spectrum\strut
\end{minipage}\tabularnewline
\begin{minipage}[t]{0.30\columnwidth}\raggedright
\texttt{BIN\_HZ}\strut
\end{minipage} &
\begin{minipage}[t]{0.30\columnwidth}\raggedright
10.766\ldots{}\strut
\end{minipage} &
\begin{minipage}[t]{0.30\columnwidth}\raggedright
Hz per bin (\texttt{SAMPLE\_RATE\ /\ NFFT})\strut
\end{minipage}\tabularnewline
\begin{minipage}[t]{0.30\columnwidth}\raggedright
\texttt{SHAVE}\strut
\end{minipage} &
\begin{minipage}[t]{0.30\columnwidth}\raggedright
0.15\strut
\end{minipage} &
\begin{minipage}[t]{0.30\columnwidth}\raggedright
v1 (and fallback): fraction trimmed from each end of a
note\strut
\end{minipage}\tabularnewline
\begin{minipage}[t]{0.30\columnwidth}\raggedright
\texttt{SMOOTH\_FRAMES}\strut
\end{minipage} &
\begin{minipage}[t]{0.30\columnwidth}\raggedright
17\strut
\end{minipage} &
\begin{minipage}[t]{0.30\columnwidth}\raggedright
v2: level smoothed over 17 frames ($\approx$ 0.2 s) before finding
the sustain\strut
\end{minipage}\tabularnewline
\begin{minipage}[t]{0.30\columnwidth}\raggedright
\texttt{LOUD\_DB}\strut
\end{minipage} &
\begin{minipage}[t]{0.30\columnwidth}\raggedright
10\strut
\end{minipage} &
\begin{minipage}[t]{0.30\columnwidth}\raggedright
v2: reference level = median of frames within 10 dB of the
peak\strut
\end{minipage}\tabularnewline
\begin{minipage}[t]{0.30\columnwidth}\raggedright
\texttt{SUSTAIN\_DB}\strut
\end{minipage} &
\begin{minipage}[t]{0.30\columnwidth}\raggedright
4\strut
\end{minipage} &
\begin{minipage}[t]{0.30\columnwidth}\raggedright
v2: sustain = frames within ±4 dB of the reference\strut
\end{minipage}\tabularnewline
\begin{minipage}[t]{0.30\columnwidth}\raggedright
\texttt{EDGE\_FRAMES}\strut
\end{minipage} &
\begin{minipage}[t]{0.30\columnwidth}\raggedright
2\strut
\end{minipage} &
\begin{minipage}[t]{0.30\columnwidth}\raggedright
v2: frames dropped at each end of the sustain\strut
\end{minipage}\tabularnewline
\begin{minipage}[t]{0.30\columnwidth}\raggedright
\texttt{MIN\_FRAMES}\strut
\end{minipage} &
\begin{minipage}[t]{0.30\columnwidth}\raggedright
3\strut
\end{minipage} &
\begin{minipage}[t]{0.30\columnwidth}\raggedright
v2: fewer frames than this $\rightarrow$ fixed-trim fallback\strut
\end{minipage}\tabularnewline
\begin{minipage}[t]{0.30\columnwidth}\raggedright
\texttt{DECAY\_DB}, \texttt{DECAY\_MAX\_S}\strut
\end{minipage} &
\begin{minipage}[t]{0.30\columnwidth}\raggedright
20, 0.5\strut
\end{minipage} &
\begin{minipage}[t]{0.30\columnwidth}\raggedright
v2, struck sounds: window from the peak until 20 dB down, at
most 0.5 s\strut
\end{minipage}\tabularnewline
\bottomrule
\end{longtable}

\texttt{LABORCH\_ROOT} is read once, at import time.

\subsubsection{\texorpdfstring{5.2
\texttt{laborch.bank}}{5.2 laborch.bank}}

\paragraph{Constants}

\textbf{\texttt{INSTRUMENT\_NAMES}} ---
\texttt{dict{[}str,\ str{]}}, the 33 instrument codes of the
current Bank mapped to names,
e.g.~\texttt{"vln"}~$\rightarrow$~\texttt{"Violin"},
\texttt{"clb"}~$\rightarrow$~\texttt{"Bb\ Clarinet"},
\texttt{"tpc"}~$\rightarrow$~\texttt{"C\ Trumpet"}.

\textbf{\texttt{MODE\_NAMES}} ---
\texttt{dict{[}str,\ str{]}}, performance-mode codes:
\texttt{ord} ordinary (winds: with vibrato), \texttt{mtd}
muted, \texttt{nhr} natural harmonic, \texttt{ahr}
artificial harmonic, \texttt{hmo} harmon mute (open),
\texttt{stm} straight mute, \texttt{ped} pedal tone,
\texttt{bwd} bowed, \texttt{hdm} hard mallet, \texttt{bbt}
brass beater. Confirmed against François' Instrument Index
spreadsheet.

\textbf{\texttt{VARIANT\_NAMES}} --- \texttt{o} open string,
\texttt{x} stopped (the same pitch fingered on the next
string down). Variants occur only on open-string pitches.

\paragraph{\texorpdfstring{\texttt{parse\_name(name)\ -\textgreater{}\ SoundName}}{parse\_name(name) -\textgreater{} SoundName}}

Parses a Bank sound name, with or without extension.

\begin{Shaded}
\begin{Highlighting}[]
\OperatorTok{\textgreater{}\textgreater{}\textgreater{}}\NormalTok{ s }\OperatorTok{=}\NormalTok{ parse\_name(}\StringTok{"ctb{-}2fs4{-}nhr{-}mf.txt"}\NormalTok{)}
\OperatorTok{\textgreater{}\textgreater{}\textgreater{}}\NormalTok{ s.instrument, s.string, s.pitch, s.mode, s.dynamic}
\NormalTok{(}\StringTok{\textquotesingle{}ctb\textquotesingle{}}\NormalTok{, }\DecValTok{2}\NormalTok{, }\StringTok{\textquotesingle{}fs4\textquotesingle{}}\NormalTok{, }\StringTok{\textquotesingle{}nhr\textquotesingle{}}\NormalTok{, }\StringTok{\textquotesingle{}mf\textquotesingle{}}\NormalTok{)}
\end{Highlighting}
\end{Shaded}

Raises \texttt{ValueError} for names that do not follow the
scheme.

\paragraph{\texorpdfstring{\texttt{class\ SoundName} (frozen
dataclass)}{class SoundName (frozen dataclass)}}

\begin{longtable}[]{@{}lll@{}}
\toprule
\begin{minipage}[b]{0.30\columnwidth}\raggedright
Field\strut
\end{minipage} &
\begin{minipage}[b]{0.30\columnwidth}\raggedright
Type\strut
\end{minipage} &
\begin{minipage}[b]{0.30\columnwidth}\raggedright
Example\strut
\end{minipage}\tabularnewline
\midrule
\endhead
\begin{minipage}[t]{0.30\columnwidth}\raggedright
\texttt{instrument}\strut
\end{minipage} &
\begin{minipage}[t]{0.30\columnwidth}\raggedright
\texttt{str}\strut
\end{minipage} &
\begin{minipage}[t]{0.30\columnwidth}\raggedright
\texttt{"vln"}\strut
\end{minipage}\tabularnewline
\begin{minipage}[t]{0.30\columnwidth}\raggedright
\texttt{note}\strut
\end{minipage} &
\begin{minipage}[t]{0.30\columnwidth}\raggedright
\texttt{str}\strut
\end{minipage} &
\begin{minipage}[t]{0.30\columnwidth}\raggedright
\texttt{"a"}, \texttt{"cs"} (sharps as \texttt{s})\strut
\end{minipage}\tabularnewline
\begin{minipage}[t]{0.30\columnwidth}\raggedright
\texttt{octave}\strut
\end{minipage} &
\begin{minipage}[t]{0.30\columnwidth}\raggedright
\texttt{int}\strut
\end{minipage} &
\begin{minipage}[t]{0.30\columnwidth}\raggedright
\texttt{4}\strut
\end{minipage}\tabularnewline
\begin{minipage}[t]{0.30\columnwidth}\raggedright
\texttt{mode}\strut
\end{minipage} &
\begin{minipage}[t]{0.30\columnwidth}\raggedright
\texttt{str}\strut
\end{minipage} &
\begin{minipage}[t]{0.30\columnwidth}\raggedright
\texttt{"ord"}\strut
\end{minipage}\tabularnewline
\begin{minipage}[t]{0.30\columnwidth}\raggedright
\texttt{dynamic}\strut
\end{minipage} &
\begin{minipage}[t]{0.30\columnwidth}\raggedright
\texttt{str}\strut
\end{minipage} &
\begin{minipage}[t]{0.30\columnwidth}\raggedright
\texttt{"pp"}, \texttt{"mf"}, \texttt{"ff"}\strut
\end{minipage}\tabularnewline
\begin{minipage}[t]{0.30\columnwidth}\raggedright
\texttt{string}\strut
\end{minipage} &
\begin{minipage}[t]{0.30\columnwidth}\raggedright
\texttt{int} or \texttt{None}\strut
\end{minipage} &
\begin{minipage}[t]{0.30\columnwidth}\raggedright
\texttt{2} in \texttt{ctb-2fs4-nhr-mf}\strut
\end{minipage}\tabularnewline
\begin{minipage}[t]{0.30\columnwidth}\raggedright
\texttt{variant}\strut
\end{minipage} &
\begin{minipage}[t]{0.30\columnwidth}\raggedright
\texttt{str}\strut
\end{minipage} &
\begin{minipage}[t]{0.30\columnwidth}\raggedright
\texttt{"o"} (open string) / \texttt{"x"} (stopped) in
\texttt{vln-a4o-ord-mf}\strut
\end{minipage}\tabularnewline
\bottomrule
\end{longtable}

\begin{longtable}[]{@{}ll@{}}
\toprule
Property & Returns\tabularnewline
\midrule
\endhead
\texttt{pitch} & \texttt{"a4"}\tabularnewline
\texttt{midi} & MIDI number, C4 = 60, A4 = 69\tabularnewline
\texttt{freq} & equal-tempered fundamental in Hz, A4 =
440\tabularnewline
\texttt{str(s)} & the original name\tabularnewline
\bottomrule
\end{longtable}

\paragraph{\texorpdfstring{\texttt{class\ Palette}
(dataclass)}{class Palette (dataclass)}}

\begin{longtable}[]{@{}lll@{}}
\toprule
\begin{minipage}[b]{0.30\columnwidth}\raggedright
Attribute\strut
\end{minipage} &
\begin{minipage}[b]{0.30\columnwidth}\raggedright
Type\strut
\end{minipage} &
\begin{minipage}[b]{0.30\columnwidth}\raggedright
Meaning\strut
\end{minipage}\tabularnewline
\midrule
\endhead
\begin{minipage}[t]{0.30\columnwidth}\raggedright
\texttt{names}\strut
\end{minipage} &
\begin{minipage}[t]{0.30\columnwidth}\raggedright
\texttt{list{[}str{]}}\strut
\end{minipage} &
\begin{minipage}[t]{0.30\columnwidth}\raggedright
sound names, one per column\strut
\end{minipage}\tabularnewline
\begin{minipage}[t]{0.30\columnwidth}\raggedright
\texttt{matrix}\strut
\end{minipage} &
\begin{minipage}[t]{0.30\columnwidth}\raggedright
\texttt{ndarray\ (2048,\ n)}\strut
\end{minipage} &
\begin{minipage}[t]{0.30\columnwidth}\raggedright
spectra as \textbf{columns} (the basis X of RH2009)\strut
\end{minipage}\tabularnewline
\bottomrule
\end{longtable}

\texttt{len(palette)} is the number of sounds. Palettes are
normally made with \texttt{Bank.select}, but any set of
spectra can be wrapped:
\texttt{Palette({[}"a",\ "b"{]},\ np.column\_stack({[}spec\_a,\ spec\_b{]}))}.

\paragraph{\texorpdfstring{\texttt{class\ Bank}}{class Bank}}

All Bank spectra, one \textbf{row} per sound, rows sorted by
name.

\textbf{Constructor}
\texttt{Bank(names,\ spectra,\ source="",\ kind="magnitude",\ version=1,\ meta=None)}
--- \texttt{spectra} must have shape
\texttt{(len(names),\ 2048)}; every name must parse with
\texttt{parse\_name}; \texttt{kind} is \texttt{"power"} or
\texttt{"magnitude"}.

\textbf{Class methods}

\begin{longtable}[]{@{}ll@{}}
\toprule
\begin{minipage}[b]{0.47\columnwidth}\raggedright
Method\strut
\end{minipage} &
\begin{minipage}[b]{0.47\columnwidth}\raggedright
Description\strut
\end{minipage}\tabularnewline
\midrule
\endhead
\begin{minipage}[t]{0.47\columnwidth}\raggedright
\texttt{Bank.load(version=2,\ normalize=None,\ **kw)}\strut
\end{minipage} &
\begin{minipage}[t]{0.47\columnwidth}\raggedright
Loads Bank v2 (default) or v1. \texttt{normalize="unit"}
divides every spectrum by its sum, keeping only its shape.
Other keywords go to \texttt{load\_v2} /
\texttt{load\_v1}.\strut
\end{minipage}\tabularnewline
\begin{minipage}[t]{0.47\columnwidth}\raggedright
\texttt{Bank.load\_v2(path=BANK\_FILE,\ meta\_path=BANK\_META)}\strut
\end{minipage} &
\begin{minipage}[t]{0.47\columnwidth}\raggedright
Bank v2: power spectra at recorded level
(\texttt{kind="power"}) and the metadata table. Raises
\texttt{FileNotFoundError} if \texttt{bank\_v2.npz} is
missing (build it with
\texttt{python\ -m\ laborch.build}).\strut
\end{minipage}\tabularnewline
\begin{minipage}[t]{0.47\columnwidth}\raggedright
\texttt{Bank.load\_v1(fft\_dir=FFT\_DIR,\ use\_cache=True,\ cache\_file=CACHE\_FILE)}\strut
\end{minipage} &
\begin{minipage}[t]{0.47\columnwidth}\raggedright
Bank v1 from \texttt{FFTBank/*.txt} (unit-sum magnitudes).
The first call reads 1,515 text files (about a minute on
Dropbox) and writes \texttt{.cache/fftbank.npz}; later calls
load the cache in about a second.\strut
\end{minipage}\tabularnewline
\begin{minipage}[t]{0.47\columnwidth}\raggedright
\texttt{Bank.load\_old(nbb\_dir=OLD\_FFT\_DIR)}\strut
\end{minipage} &
\begin{minipage}[t]{0.47\columnwidth}\raggedright
Loads the 2009 \texttt{.nbb} spectra (2,048 float32 values
each). One malformed name (\texttt{obo-f5-c6.nbb}) is
skipped with a warning. Old instrument codes differ
(\texttt{cla}, \texttt{hrn}, \texttt{tpt}, \texttt{mar},
\texttt{btr}).\strut
\end{minipage}\tabularnewline
\bottomrule
\end{longtable}

\textbf{Attributes}

\begin{longtable}[]{@{}ll@{}}
\toprule
\begin{minipage}[b]{0.47\columnwidth}\raggedright
Attribute\strut
\end{minipage} &
\begin{minipage}[b]{0.47\columnwidth}\raggedright
Meaning\strut
\end{minipage}\tabularnewline
\midrule
\endhead
\begin{minipage}[t]{0.47\columnwidth}\raggedright
\texttt{names}\strut
\end{minipage} &
\begin{minipage}[t]{0.47\columnwidth}\raggedright
sorted list of sound names\strut
\end{minipage}\tabularnewline
\begin{minipage}[t]{0.47\columnwidth}\raggedright
\texttt{spectra}\strut
\end{minipage} &
\begin{minipage}[t]{0.47\columnwidth}\raggedright
\texttt{ndarray\ (n\_sounds,\ 2048)}\strut
\end{minipage}\tabularnewline
\begin{minipage}[t]{0.47\columnwidth}\raggedright
\texttt{info}\strut
\end{minipage} &
\begin{minipage}[t]{0.47\columnwidth}\raggedright
\texttt{dict{[}name,\ SoundName{]}}\strut
\end{minipage}\tabularnewline
\begin{minipage}[t]{0.47\columnwidth}\raggedright
\texttt{source}\strut
\end{minipage} &
\begin{minipage}[t]{0.47\columnwidth}\raggedright
file or folder the Bank was read from\strut
\end{minipage}\tabularnewline
\begin{minipage}[t]{0.47\columnwidth}\raggedright
\texttt{kind}\strut
\end{minipage} &
\begin{minipage}[t]{0.47\columnwidth}\raggedright
\texttt{"power"} (v2) or \texttt{"magnitude"} (v1,
2009)\strut
\end{minipage}\tabularnewline
\begin{minipage}[t]{0.47\columnwidth}\raggedright
\texttt{version}\strut
\end{minipage} &
\begin{minipage}[t]{0.47\columnwidth}\raggedright
2 or 1\strut
\end{minipage}\tabularnewline
\begin{minipage}[t]{0.47\columnwidth}\raggedright
\texttt{meta}\strut
\end{minipage} &
\begin{minipage}[t]{0.47\columnwidth}\raggedright
\texttt{dict{[}name,\ dict{]}}, the metadata row of each
sound (v2; empty for v1)\strut
\end{minipage}\tabularnewline
\begin{minipage}[t]{0.47\columnwidth}\raggedright
\texttt{instruments}\strut
\end{minipage} &
\begin{minipage}[t]{0.47\columnwidth}\raggedright
sorted list of instrument codes present (property)\strut
\end{minipage}\tabularnewline
\bottomrule
\end{longtable}

\textbf{Methods}

\begin{longtable}[]{@{}lll@{}}
\toprule
\begin{minipage}[b]{0.30\columnwidth}\raggedright
Method\strut
\end{minipage} &
\begin{minipage}[b]{0.30\columnwidth}\raggedright
Returns\strut
\end{minipage} &
\begin{minipage}[b]{0.30\columnwidth}\raggedright
Description\strut
\end{minipage}\tabularnewline
\midrule
\endhead
\begin{minipage}[t]{0.30\columnwidth}\raggedright
\texttt{len(bank)}\strut
\end{minipage} &
\begin{minipage}[t]{0.30\columnwidth}\raggedright
\texttt{int}\strut
\end{minipage} &
\begin{minipage}[t]{0.30\columnwidth}\raggedright
number of sounds\strut
\end{minipage}\tabularnewline
\begin{minipage}[t]{0.30\columnwidth}\raggedright
\texttt{name\ in\ bank}\strut
\end{minipage} &
\begin{minipage}[t]{0.30\columnwidth}\raggedright
\texttt{bool}\strut
\end{minipage} &
\begin{minipage}[t]{0.30\columnwidth}\raggedright
membership\strut
\end{minipage}\tabularnewline
\begin{minipage}[t]{0.30\columnwidth}\raggedright
\texttt{spectrum(name)}\strut
\end{minipage} &
\begin{minipage}[t]{0.30\columnwidth}\raggedright
\texttt{ndarray\ (2048,)}\strut
\end{minipage} &
\begin{minipage}[t]{0.30\columnwidth}\raggedright
one sound's spectrum (a view; copy before modifying)\strut
\end{minipage}\tabularnewline
\begin{minipage}[t]{0.30\columnwidth}\raggedright
\texttt{select(instruments=None,\ modes=None,\ dynamics=None,\ midi\_range=None,\ exclude=(),\ families=None)}\strut
\end{minipage} &
\begin{minipage}[t]{0.30\columnwidth}\raggedright
\texttt{Palette}\strut
\end{minipage} &
\begin{minipage}[t]{0.30\columnwidth}\raggedright
builds a Palette. Each filter is optional and takes a list
or a single string; \texttt{midi\_range=(lo,\ hi)} is
inclusive; \texttt{exclude} removes named sounds, e.g.~the
Target; \texttt{families} selects by family (see 5.7) and
combines with \texttt{instruments} as an intersection.
Raises \texttt{ValueError} if nothing matches.\strut
\end{minipage}\tabularnewline
\begin{minipage}[t]{0.30\columnwidth}\raggedright
\texttt{level\_db(name)}\strut
\end{minipage} &
\begin{minipage}[t]{0.30\columnwidth}\raggedright
\texttt{float}\strut
\end{minipage} &
\begin{minipage}[t]{0.30\columnwidth}\raggedright
10 log₁₀ of the spectrum's sum: the sound's level in dBFS
for v2\strut
\end{minipage}\tabularnewline
\begin{minipage}[t]{0.30\columnwidth}\raggedright
\texttt{sound\_path(name)}\strut
\end{minipage} &
\begin{minipage}[t]{0.30\columnwidth}\raggedright
\texttt{Path}\strut
\end{minipage} &
\begin{minipage}[t]{0.30\columnwidth}\raggedright
the \texttt{.aif} recording for a sound;
\texttt{FileNotFoundError} if absent\strut
\end{minipage}\tabularnewline
\bottomrule
\end{longtable}

\begin{Shaded}
\begin{Highlighting}[]
\NormalTok{quartet }\OperatorTok{=}\NormalTok{ bank.select(instruments}\OperatorTok{=}\NormalTok{[}\StringTok{"vln"}\NormalTok{, }\StringTok{"vla"}\NormalTok{, }\StringTok{"vlc"}\NormalTok{],}
\NormalTok{                      modes}\OperatorTok{=}\NormalTok{[}\StringTok{"ord"}\NormalTok{, }\StringTok{"mtd"}\NormalTok{],}
\NormalTok{                      midi\_range}\OperatorTok{=}\NormalTok{(}\DecValTok{48}\NormalTok{, }\DecValTok{84}\NormalTok{))}
\end{Highlighting}
\end{Shaded}

\paragraph{\texorpdfstring{\texttt{frequencies()\ -\textgreater{}\ ndarray}}{frequencies() -\textgreater{} ndarray}}

Centre frequency in Hz of each of the 2,048 bins (bin 0 =
DC). Use it as the x-axis when plotting spectra.

\subsubsection{\texorpdfstring{5.3
\texttt{laborch.spectra}}{5.3 laborch.spectra}}

The averaged-spectrum pipelines used to build the Banks,
plus audio I/O. Functions marked \textdagger{} need pydub and ffmpeg.

\paragraph{\texorpdfstring{Bank v2:
\texttt{power\_spectrum(x,\ envelope="sustained",\ region=None)\ -\textgreater{}\ (ndarray,\ Sustain)}}{Bank v2: power\_spectrum(x, envelope="sustained", region=None) -\textgreater{} (ndarray, Sustain)}}

The Bank v2 recipe:

\begin{enumerate}
\def\labelenumi{\arabic{enumi}.}
\tightlist
\item
  cut \texttt{x} (mono, scaled to −1..1; see
  \texttt{load\_float}) into 4,096-sample Hamming frames at
  hop 512;
\item
  take each frame's one-sided power spectrum, scaled so it
  sums to the frame's mean square (Parseval);
\item
  choose the frames to use with \texttt{find\_sustain} (or
  the \texttt{region} \texttt{(start\_s,\ end\_s)} you
  give);
\item
  average the frame powers.
\end{enumerate}

\texttt{10*log10(spec.sum())} is the level of the region in
dBFS, and spectra of uncorrelated sounds add. Returns the
spectrum and a \texttt{Sustain}.

\begin{Shaded}
\begin{Highlighting}[]
\ImportTok{from}\NormalTok{ laborch.spectra }\ImportTok{import}\NormalTok{ power\_of\_file}
\NormalTok{target, sus }\OperatorTok{=}\NormalTok{ power\_of\_file(}\StringTok{"my\_target.wav"}\NormalTok{)      }\CommentTok{\# a Target in Bank v2 units}
\NormalTok{sus.start\_s, sus.end\_s, sus.level\_db, sus.method}
\end{Highlighting}
\end{Shaded}

\paragraph{\texorpdfstring{\texttt{find\_sustain(ms,\ envelope="sustained")\ -\textgreater{}\ Sustain}}{find\_sustain(ms, envelope="sustained") -\textgreater{} Sustain}}

Chooses frames from their mean-square levels \texttt{ms}
(from \texttt{frame\_power}).

\begin{itemize}
\tightlist
\item
  \texttt{"sustained"} (bowed, blown): smooth the dB level
  over \texttt{SMOOTH\_FRAMES} ($\approx$ 0.2 s, so vibrato and
  tremolo do not break it up); take as reference the median
  of frames within \texttt{LOUD\_DB} of the smoothed peak;
  keep the longest run within ±\texttt{SUSTAIN\_DB} of it,
  less \texttt{EDGE\_FRAMES} at each end. Fewer than
  \texttt{MIN\_FRAMES} $\rightarrow$ the v1 fixed 15\% trim
  (\texttt{method="trim"}).
\item
  \texttt{"decaying"} (struck): from the loudest frame until
  \texttt{DECAY\_DB} below it, at most
  \texttt{DECAY\_MAX\_S}.
\end{itemize}

Which envelope a Bank sound has is in its metadata
(\texttt{metadata.envelope\_class}).

\paragraph{\texorpdfstring{\texttt{class\ Sustain}
(dataclass)}{class Sustain (dataclass)}}

\texttt{first}, \texttt{last} (frame indices, inclusive),
\texttt{method} (\texttt{"sustain"}, \texttt{"decay"},
\texttt{"trim"}, \texttt{"given"}), \texttt{level\_db}
(mean-square level of the region, dBFS), \texttt{peak\_db}
(loudest frame); properties \texttt{n\_frames},
\texttt{start\_s}, \texttt{end\_s} (from the start of the
first window to the end of the last).

\paragraph{\texorpdfstring{\texttt{power\_of\_file(path,\ envelope="sustained",\ region=None)}
†, \texttt{load\_float(path)} †,
\texttt{segment\_float(segment)}~\textsuperscript{\textdagger}}{power\_of\_file(path, envelope="sustained", region=None) †, load\_float(path) †, segment\_float(segment) (needs pydub)}}

\texttt{power\_spectrum} of a file; mono samples scaled to
−1..1.

\paragraph{\texorpdfstring{\texttt{frame\_power(x)\ -\textgreater{}\ (P,\ ms)}
and
\texttt{frames(x)}}{frame\_power(x) -\textgreater{} (P, ms) and frames(x)}}

Per-frame power spectra \texttt{(n\_frames,\ 2048)} and mean
squares; the frame matrix.

\paragraph{Bank v1}

\paragraph{\texorpdfstring{\texttt{make\_fft(samples,\ shave=0.15,\ normalize=True)\ -\textgreater{}\ ndarray}}{make\_fft(samples, shave=0.15, normalize=True) -\textgreater{} ndarray}}

The Bank's spectrum recipe (notebook 1.11):

\begin{enumerate}
\def\labelenumi{\arabic{enumi}.}
\tightlist
\item
  drop the first and last \texttt{shave} fraction of the
  samples (attack and decay);
\item
  slide a 4,096-sample Hamming window in steps of 512;
\item
  combine the frames as $\bigl(\sum_{\text{frames}} |X_k|^2\bigr)^{1/2} / n$,
  keeping bins 0--2047;
\item
  if \texttt{normalize}, divide by the sum (unit area), as
  every Bank file is.
\end{enumerate}

\texttt{samples} is a 1-D array of mono samples at 44.1 kHz.
Raises \texttt{ValueError} if the sound is shorter than one
window after trimming. Use \texttt{normalize=False} to keep
the absolute level.

\paragraph{\texorpdfstring{\texttt{fft\_of\_file(path,\ **kw)\ -\textgreater{}\ ndarray}~\textsuperscript{\textdagger}}{fft\_of\_file(path, **kw) -\textgreater{} ndarray (needs pydub)}}

\texttt{make\_fft(load\_samples(path),\ **kw)}. Recomputing
a Bank file this way matches the stored spectrum to about
0.1\% of its peak.

\paragraph{\texorpdfstring{\texttt{load\_segment(path)\ -\textgreater{}\ pydub.AudioSegment}~\textsuperscript{\textdagger}}{load\_segment(path) -\textgreater{} pydub.AudioSegment (needs pydub)}}

Reads any ffmpeg-readable file (AIF, WAV, MP3, \ldots).

\paragraph{\texorpdfstring{\texttt{segment\_samples(segment)\ -\textgreater{}\ ndarray}~\textsuperscript{\textdagger}}{segment\_samples(segment) -\textgreater{} ndarray (needs pydub)}}

Mono samples of an \texttt{AudioSegment} as floats in the
file's integer scale (the scale used to build the Bank).
Warns if the sample rate is not 44.1 kHz.

\paragraph{\texorpdfstring{\texttt{load\_samples(path)\ -\textgreater{}\ ndarray}~\textsuperscript{\textdagger}}{load\_samples(path) -\textgreater{} ndarray (needs pydub)}}

\texttt{segment\_samples(load\_segment(path))}.

\paragraph{\texorpdfstring{\texttt{split\_on\_silence(path,\ min\_silence\_len=500,\ silence\_thresh=-50,\ keep\_silence=100)\ -\textgreater{}\ list}~\textsuperscript{\textdagger}}{split\_on\_silence(path, min\_silence\_len=500, silence\_thresh=-50, keep\_silence=100) -\textgreater{} list (needs pydub)}}

Splits a multi-note recording into one \texttt{AudioSegment}
per note (notebooks 1.8, 1.9, 1.12). Times in ms, threshold
in dBFS.

\begin{Shaded}
\begin{Highlighting}[]
\ImportTok{from}\NormalTok{ laborch.spectra }\ImportTok{import}\NormalTok{ split\_on\_silence, segment\_samples, make\_fft}
\NormalTok{notes }\OperatorTok{=}\NormalTok{ split\_on\_silence(}\StringTok{"MiscSounds/Trumpet{-}mutes\_Basson{-}multiphics/11 Track 11.aiff"}\NormalTok{)}
\NormalTok{cup\_c5 }\OperatorTok{=}\NormalTok{ make\_fft(segment\_samples(notes[}\DecValTok{2}\NormalTok{]))}
\end{Highlighting}
\end{Shaded}

\paragraph{\texorpdfstring{\texttt{spectral\_peaks(spec,\ rel\_height=0.02,\ max\_hz=None)\ -\textgreater{}\ list}}{spectral\_peaks(spec, rel\_height=0.02, max\_hz=None) -\textgreater{} list}}

Local maxima of a spectrum as
\texttt{{[}(Hz,\ height\ relative\ to\ the\ tallest),\ ...{]}},
lowest frequency first. Peaks below \texttt{rel\_height} $\times$
the tallest, or above \texttt{max\_hz}, are dropped.

\begin{Shaded}
\begin{Highlighting}[]
\OperatorTok{\textgreater{}\textgreater{}\textgreater{}}\NormalTok{ spectral\_peaks(make\_fft(tone(}\FloatTok{130.81}\NormalTok{, [}\DecValTok{1}\NormalTok{, }\DecValTok{4}\NormalTok{, }\DecValTok{10}\NormalTok{], [}\DecValTok{1}\NormalTok{, }\FloatTok{.5}\NormalTok{, }\FloatTok{.25}\NormalTok{])), max\_hz}\OperatorTok{=}\DecValTok{3000}\NormalTok{)}
\NormalTok{[(}\FloatTok{129.2}\NormalTok{, }\FloatTok{1.0}\NormalTok{), (}\FloatTok{527.6}\NormalTok{, }\FloatTok{0.45}\NormalTok{), (}\FloatTok{1302.8}\NormalTok{, }\FloatTok{0.21}\NormalTok{)]}
\end{Highlighting}
\end{Shaded}

\subsubsection{\texorpdfstring{5.4
\texttt{laborch.solve}}{5.4 laborch.solve}}

\paragraph{\texorpdfstring{\texttt{nnls\_fit(palette,\ target,\ sigma=None,\ maxiter=None)\ -\textgreater{}\ Fit}}{nnls\_fit(palette, target, sigma=None, maxiter=None) -\textgreater{} Fit}}

Non-negative least squares (Lawson--Hanson, via
\texttt{scipy.optimize.nnls}):

\[
  \min_{\alpha}\; \sum_i \left( \frac{T_i - \sum_j X_{ij}\,\alpha_j}{\sigma_i} \right)^{2}
  \quad \text{subject to} \quad \alpha_j \ge 0
\]

\texttt{sigma} is an optional array of 2,048 positive
per-bin uncertainties; \texttt{None} weights every bin
equally (the 2020 behaviour). This is the main solver.

\paragraph{\texorpdfstring{\texttt{pinv\_fit(palette,\ target,\ sigma=None)\ -\textgreater{}\ Fit}}{pinv\_fit(palette, target, sigma=None) -\textgreater{} Fit}}

Unconstrained least squares via the pseudo-inverse ---
RH2009's ``SVD'' method. Weights can be negative (an
instrument ``subtracting'' sound), so use it only for
comparison.

\paragraph{\texorpdfstring{\texttt{dot(palette,\ target,\ normalize=False)\ -\textgreater{}\ Fit}}{dot(palette, target, normalize=False) -\textgreater{} Fit}}

RH2009's ``DOT'' method: the projection \mbox{$p_j = \sum_i T_i X_{ij}$} of
the Target on each sound alone. With \texttt{normalize=True}
it is cosine similarity (shape only). The weights are
scores, not mixing amounts; use DOT to rank or prefilter,
not to build a mix.

\paragraph{\texorpdfstring{\texttt{chi2(target,\ model,\ sigma=None)\ -\textgreater{}\ float}}{chi2(target, model, sigma=None) -\textgreater{} float}}

Sum of squared, optionally $\sigma$-weighted, residuals.

\paragraph{\texorpdfstring{\texttt{class\ Fit}
(dataclass)}{class Fit (dataclass)}}

The result of any solver: a Recommendation before
playability rules.

\begin{longtable}[]{@{}ll@{}}
\toprule
Attribute & Meaning\tabularnewline
\midrule
\endhead
\texttt{names} & Palette sound names\tabularnewline
\texttt{weights} & \texttt{ndarray}, one weight per name
(same order)\tabularnewline
\texttt{target} & the Target spectrum\tabularnewline
\texttt{matrix} & the Palette matrix used\tabularnewline
\texttt{method} & \texttt{"nnls"}, \texttt{"pinv"},
\texttt{"dot"}, \texttt{"dot-cos"}, or \texttt{"nnls/top3"}
etc.\tabularnewline
\texttt{sigma} & the $\sigma$ used, or \texttt{None}\tabularnewline
\bottomrule
\end{longtable}

\begin{longtable}[]{@{}ll@{}}
\toprule
Property & Returns\tabularnewline
\midrule
\endhead
\texttt{model} & fitted spectrum
\texttt{matrix\ @\ weights}\tabularnewline
\texttt{chi2} & $\chi^2$ of the fit ($\sigma$-weighted if $\sigma$ was
given)\tabularnewline
\texttt{relative\_error} & $\lVert T - \text{model} \rVert / \lVert T \rVert$, unweighted; 0
is perfect\tabularnewline
\texttt{order} & indices sorted by weight, largest
first\tabularnewline
\bottomrule
\end{longtable}

\begin{longtable}[]{@{}lll@{}}
\toprule
\begin{minipage}[b]{0.30\columnwidth}\raggedright
Method\strut
\end{minipage} &
\begin{minipage}[b]{0.30\columnwidth}\raggedright
Returns\strut
\end{minipage} &
\begin{minipage}[b]{0.30\columnwidth}\raggedright
Description\strut
\end{minipage}\tabularnewline
\midrule
\endhead
\begin{minipage}[t]{0.30\columnwidth}\raggedright
\texttt{ranked(n=None,\ nonzero=True,\ rtol=0.0)}\strut
\end{minipage} &
\begin{minipage}[t]{0.30\columnwidth}\raggedright
\texttt{list{[}(name,\ weight){]}}\strut
\end{minipage} &
\begin{minipage}[t]{0.30\columnwidth}\raggedright
largest first; with \texttt{nonzero}, drops weights $\le$
\texttt{rtol}~$\times$~\texttt{max} (use \texttt{rtol=1e-9} to hide
round-off)\strut
\end{minipage}\tabularnewline
\begin{minipage}[t]{0.30\columnwidth}\raggedright
\texttt{top(n,\ refit=False)}\strut
\end{minipage} &
\begin{minipage}[t]{0.30\columnwidth}\raggedright
\texttt{Fit}\strut
\end{minipage} &
\begin{minipage}[t]{0.30\columnwidth}\raggedright
keeps the n largest weights. \texttt{refit=False} keeps
their values (as the 2020 mock-up did); \texttt{refit=True}
re-solves NNLS on those n sounds, which never increases
$\chi^2$\strut
\end{minipage}\tabularnewline
\begin{minipage}[t]{0.30\columnwidth}\raggedright
\texttt{kept\_fraction(n)}\strut
\end{minipage} &
\begin{minipage}[t]{0.30\columnwidth}\raggedright
\texttt{float}\strut
\end{minipage} &
\begin{minipage}[t]{0.30\columnwidth}\raggedright
share of total weight in the top n\strut
\end{minipage}\tabularnewline
\begin{minipage}[t]{0.30\columnwidth}\raggedright
\texttt{summary(n=10)}\strut
\end{minipage} &
\begin{minipage}[t]{0.30\columnwidth}\raggedright
\texttt{str}\strut
\end{minipage} &
\begin{minipage}[t]{0.30\columnwidth}\raggedright
readable report\strut
\end{minipage}\tabularnewline
\bottomrule
\end{longtable}

\subsubsection{\texorpdfstring{5.5
\texttt{laborch.audio}}{5.5 laborch.audio}}

Needs pydub and ffmpeg (except \texttt{weights\_to\_db}).

\paragraph{\texorpdfstring{\texttt{weights\_to\_db(weights,\ floor\_db=-60.0,\ kind="magnitude",\ relative=True)\ -\textgreater{}\ ndarray}}{weights\_to\_db(weights, floor\_db=-60.0, kind="magnitude", relative=True) -\textgreater{} ndarray}}

Gain in dB of each weight: 20 log₁₀ w for magnitude spectra
(v1), 10 log₁₀ w for power spectra (v2). With
\texttt{relative=True} the largest weight is 0 dB; with
\texttt{relative=False} the gain is relative to each
recording's own level, which is what a v2 weight means. Zero
weights get \texttt{floor\_db}.

\paragraph{\texorpdfstring{\texttt{mix(bank,\ names,\ gains\_db)\ -\textgreater{}\ AudioSegment}}{mix(bank, names, gains\_db) -\textgreater{} AudioSegment}}

Overlays the named Bank recordings, each shifted by its gain
in dB. The result is as long as the longest sound.

\paragraph{\texorpdfstring{\texttt{render\_fit(bank,\ fit,\ n=8,\ path=None,\ headroom\_db=0.0)\ -\textgreater{}\ AudioSegment}}{render\_fit(bank, fit, n=8, path=None, headroom\_db=0.0) -\textgreater{} AudioSegment}}

Mixes the top-n sounds of a \texttt{Fit}; writes a file if
\texttt{path} is given (format from its extension,
e.g.~\texttt{.wav}, \texttt{.mp3}). In Jupyter, the returned
segment plays inline. With Bank v2 the gains are absolute,
10 log₁₀ w on each recording; if the loudest would exceed
\texttt{headroom\_db}, all gains are lowered together,
keeping the balance. With Bank v1 the gains are relative to
the largest weight.

\subsubsection{\texorpdfstring{5.6
\texttt{laborch.synth}}{5.6 laborch.synth}}

Synthetic sounds (notebooks 1.2 and 2.1). Pure NumPy: no
pydub needed.

\paragraph{\texorpdfstring{\texttt{tone(f0,\ partials=(1,),\ amps=None,\ seconds=1.0,\ sr=44100,\ decay=0.0,\ phases=None,\ peak=1.0)\ -\textgreater{}\ ndarray}}{tone(f0, partials=(1,), amps=None, seconds=1.0, sr=44100, decay=0.0, phases=None, peak=1.0) -\textgreater{} ndarray}}

A sum of sine partials at multiples of \texttt{f0}:

\[
  x(t) = e^{-\lambda t} \sum_k a_k \sin\!\left(2\pi\, p_k f_0\, t + \phi_k\right),
  \qquad \lambda = \texttt{decay}
\]

\begin{longtable}[]{@{}ll@{}}
\toprule
\begin{minipage}[b]{0.47\columnwidth}\raggedright
Parameter\strut
\end{minipage} &
\begin{minipage}[b]{0.47\columnwidth}\raggedright
Meaning\strut
\end{minipage}\tabularnewline
\midrule
\endhead
\begin{minipage}[t]{0.47\columnwidth}\raggedright
\texttt{f0}\strut
\end{minipage} &
\begin{minipage}[t]{0.47\columnwidth}\raggedright
fundamental in Hz;
\texttt{parse\_name("mrb-c3-ord-mf").freq} gives
130.81\strut
\end{minipage}\tabularnewline
\begin{minipage}[t]{0.47\columnwidth}\raggedright
\texttt{partials}\strut
\end{minipage} &
\begin{minipage}[t]{0.47\columnwidth}\raggedright
multiples of \texttt{f0}; need not be integers (e.g.~3.9 for
a real bar)\strut
\end{minipage}\tabularnewline
\begin{minipage}[t]{0.47\columnwidth}\raggedright
\texttt{amps}\strut
\end{minipage} &
\begin{minipage}[t]{0.47\columnwidth}\raggedright
amplitude of each partial (default all 1)\strut
\end{minipage}\tabularnewline
\begin{minipage}[t]{0.47\columnwidth}\raggedright
\texttt{seconds}, \texttt{sr}\strut
\end{minipage} &
\begin{minipage}[t]{0.47\columnwidth}\raggedright
duration and sample rate\strut
\end{minipage}\tabularnewline
\begin{minipage}[t]{0.47\columnwidth}\raggedright
\texttt{decay}\strut
\end{minipage} &
\begin{minipage}[t]{0.47\columnwidth}\raggedright
exponential decay rate in 1/s; 0 = steady\strut
\end{minipage}\tabularnewline
\begin{minipage}[t]{0.47\columnwidth}\raggedright
\texttt{phases}\strut
\end{minipage} &
\begin{minipage}[t]{0.47\columnwidth}\raggedright
starting phase of each partial in radians (default 0)\strut
\end{minipage}\tabularnewline
\begin{minipage}[t]{0.47\columnwidth}\raggedright
\texttt{peak}\strut
\end{minipage} &
\begin{minipage}[t]{0.47\columnwidth}\raggedright
largest absolute sample after scaling; \texttt{2**15\ -\ 1}
gives 16-bit scale\strut
\end{minipage}\tabularnewline
\bottomrule
\end{longtable}

Raises \texttt{ValueError} if the lists differ in length or
a partial reaches the Nyquist frequency (22,050 Hz).

\paragraph{\texorpdfstring{\texttt{spectrum\_to\_audio(spec,\ seconds=1.0,\ sr=44100,\ seed=0,\ threshold=1e-4,\ peak=1.0)\ -\textgreater{}\ ndarray}}{spectrum\_to\_audio(spec, seconds=1.0, sr=44100, seed=0, threshold=1e-4, peak=1.0) -\textgreater{} ndarray}}

Additive resynthesis of a magnitude spectrum (for a v2 power
spectrum, pass \texttt{np.sqrt(spec)}). Every bin \emph{k}
above \texttt{threshold} $\times$ the maximum becomes a sine at
$k \times 10.77$~Hz with amplitude \texttt{spec{[}k{]}} and
a random phase (\texttt{seed} makes it repeatable). Works on
Bank spectra, \texttt{Fit.model} and \texttt{make\_fft}
output.

Because the Bank stores magnitudes only, the phases are
invented: the result keeps the timbre (which partials, how
strong) but not the original waveform, attack or decay. A
round trip spectrum $\rightarrow$ audio $\rightarrow$ spectrum keeps the shape with
a cosine similarity of about 0.97.

\begin{Shaded}
\begin{Highlighting}[]
\ImportTok{from}\NormalTok{ laborch.synth }\ImportTok{import}\NormalTok{ tone, spectrum\_to\_audio}
\ImportTok{from}\NormalTok{ IPython.display }\ImportTok{import}\NormalTok{ Audio}
\NormalTok{x }\OperatorTok{=}\NormalTok{ tone(}\FloatTok{130.81}\NormalTok{, partials}\OperatorTok{=}\NormalTok{[}\DecValTok{1}\NormalTok{, }\DecValTok{4}\NormalTok{, }\DecValTok{10}\NormalTok{], amps}\OperatorTok{=}\NormalTok{[}\DecValTok{1}\NormalTok{, }\FloatTok{0.5}\NormalTok{, }\FloatTok{0.25}\NormalTok{])}
\NormalTok{Audio(spectrum\_to\_audio(make\_fft(x)), rate}\OperatorTok{=}\DecValTok{44100}\NormalTok{)}
\end{Highlighting}
\end{Shaded}

\subsubsection{\texorpdfstring{5.7 \texttt{laborch.metadata}
and
\texttt{laborch.build}}{5.7 laborch.metadata and laborch.build}}

\paragraph{The metadata table}

\texttt{Soundbank\_McGill/Bank\_v2/bank\_v2.csv}, one row
per sound, also available as \texttt{bank.meta{[}name{]}}.
Readable without laborch (spreadsheet, pandas).

\begin{longtable}[]{@{}ll@{}}
\toprule
\begin{minipage}[b]{0.47\columnwidth}\raggedright
Column\strut
\end{minipage} &
\begin{minipage}[b]{0.47\columnwidth}\raggedright
Meaning\strut
\end{minipage}\tabularnewline
\midrule
\endhead
\begin{minipage}[t]{0.47\columnwidth}\raggedright
\texttt{name}, \texttt{instrument},
\texttt{instrument\_name}, \texttt{family}\strut
\end{minipage} &
\begin{minipage}[t]{0.47\columnwidth}\raggedright
sound and instrument; family from \texttt{FAMILIES}\strut
\end{minipage}\tabularnewline
\begin{minipage}[t]{0.47\columnwidth}\raggedright
\texttt{envelope}\strut
\end{minipage} &
\begin{minipage}[t]{0.47\columnwidth}\raggedright
\texttt{sustained} or \texttt{decaying} (struck: piano,
marimba, xylophone, glockenspiel, crotales, tubular bells,
hard-mallet vibraphone)\strut
\end{minipage}\tabularnewline
\begin{minipage}[t]{0.47\columnwidth}\raggedright
\texttt{note}, \texttt{octave}, \texttt{pitch},
\texttt{midi}, \texttt{freq\_hz}\strut
\end{minipage} &
\begin{minipage}[t]{0.47\columnwidth}\raggedright
pitch; \texttt{freq\_hz} equal-tempered, A4 = 440\strut
\end{minipage}\tabularnewline
\begin{minipage}[t]{0.47\columnwidth}\raggedright
\texttt{string}, \texttt{variant},
\texttt{variant\_name}\strut
\end{minipage} &
\begin{minipage}[t]{0.47\columnwidth}\raggedright
string number; \texttt{o} open / \texttt{x} stopped\strut
\end{minipage}\tabularnewline
\begin{minipage}[t]{0.47\columnwidth}\raggedright
\texttt{mode}, \texttt{mode\_name}, \texttt{dynamic}\strut
\end{minipage} &
\begin{minipage}[t]{0.47\columnwidth}\raggedright
performance mode and dynamic\strut
\end{minipage}\tabularnewline
\begin{minipage}[t]{0.47\columnwidth}\raggedright
\texttt{sub\_bank}, \texttt{aif}\strut
\end{minipage} &
\begin{minipage}[t]{0.47\columnwidth}\raggedright
McGill sub-bank and the recording's path relative to
\texttt{ROOT}\strut
\end{minipage}\tabularnewline
\begin{minipage}[t]{0.47\columnwidth}\raggedright
\texttt{duration\_s}, \texttt{sustain\_start\_s},
\texttt{sustain\_end\_s}, \texttt{n\_frames},
\texttt{method}\strut
\end{minipage} &
\begin{minipage}[t]{0.47\columnwidth}\raggedright
the note and the analysed region\strut
\end{minipage}\tabularnewline
\begin{minipage}[t]{0.47\columnwidth}\raggedright
\texttt{level\_db}, \texttt{peak\_db}\strut
\end{minipage} &
\begin{minipage}[t]{0.47\columnwidth}\raggedright
mean-square level of the region and loudest frame,
dBFS\strut
\end{minipage}\tabularnewline
\begin{minipage}[t]{0.47\columnwidth}\raggedright
\texttt{similarity\_v1}\strut
\end{minipage} &
\begin{minipage}[t]{0.47\columnwidth}\raggedright
cosine between √(v2 power) and the v1 spectrum: low = the
region changed a lot\strut
\end{minipage}\tabularnewline
\bottomrule
\end{longtable}

\paragraph{\texorpdfstring{\texttt{laborch.metadata}}{laborch.metadata}}

\begin{longtable}[]{@{}ll@{}}
\toprule
\begin{minipage}[b]{0.47\columnwidth}\raggedright
Name\strut
\end{minipage} &
\begin{minipage}[b]{0.47\columnwidth}\raggedright
Description\strut
\end{minipage}\tabularnewline
\midrule
\endhead
\begin{minipage}[t]{0.47\columnwidth}\raggedright
\texttt{FAMILIES}\strut
\end{minipage} &
\begin{minipage}[t]{0.47\columnwidth}\raggedright
\texttt{dict{[}family,\ list\ of\ codes{]}}: strings,
flutes, double reeds, clarinets, saxophones, brass,
keyboard, percussion\strut
\end{minipage}\tabularnewline
\begin{minipage}[t]{0.47\columnwidth}\raggedright
\texttt{FAMILY\_OF}\strut
\end{minipage} &
\begin{minipage}[t]{0.47\columnwidth}\raggedright
code $\rightarrow$ family\strut
\end{minipage}\tabularnewline
\begin{minipage}[t]{0.47\columnwidth}\raggedright
\texttt{envelope\_class(instrument,\ mode)}\strut
\end{minipage} &
\begin{minipage}[t]{0.47\columnwidth}\raggedright
\texttt{"sustained"} or \texttt{"decaying"}\strut
\end{minipage}\tabularnewline
\begin{minipage}[t]{0.47\columnwidth}\raggedright
\texttt{describe(name)}\strut
\end{minipage} &
\begin{minipage}[t]{0.47\columnwidth}\raggedright
the name-derived columns of a row\strut
\end{minipage}\tabularnewline
\begin{minipage}[t]{0.47\columnwidth}\raggedright
\texttt{read\_table(path=BANK\_META)},
\texttt{write\_table(rows,\ path=BANK\_META)}\strut
\end{minipage} &
\begin{minipage}[t]{0.47\columnwidth}\raggedright
CSV I/O, numbers converted\strut
\end{minipage}\tabularnewline
\begin{minipage}[t]{0.47\columnwidth}\raggedright
\texttt{ranges(rows)}\strut
\end{minipage} &
\begin{minipage}[t]{0.47\columnwidth}\raggedright
\texttt{(instrument,\ mode)\ →\ (lowest\ MIDI,\ highest\ MIDI)};
Phase 2 uses these as Palette filters\strut
\end{minipage}\tabularnewline
\begin{minipage}[t]{0.47\columnwidth}\raggedright
\texttt{spreadsheet\_names(folder)}\strut
\end{minipage} &
\begin{minipage}[t]{0.47\columnwidth}\raggedright
names listed in François' spreadsheets (needs
openpyxl)\strut
\end{minipage}\tabularnewline
\bottomrule
\end{longtable}

\begin{Shaded}
\begin{Highlighting}[]
\ImportTok{from}\NormalTok{ laborch.metadata }\ImportTok{import}\NormalTok{ ranges}
\NormalTok{ranges(bank.meta.values())[(}\StringTok{"vln"}\NormalTok{, }\StringTok{"ord"}\NormalTok{)]     }\CommentTok{\# (55, 96): g3{-}c7}
\end{Highlighting}
\end{Shaded}

\paragraph{\texorpdfstring{\texttt{laborch.build}}{laborch.build}}

\begin{Shaded}
\begin{Highlighting}[]
\ExtensionTok{python}\NormalTok{ {-}m laborch.build            \# resumable: rerun until it says it wrote the files}
\ExtensionTok{python}\NormalTok{ {-}m laborch.build {-}{-}fresh    \# start again (after changing the analysis)}
\end{Highlighting}
\end{Shaded}

Analyses every \texttt{.aif} in \texttt{SOUND\_DIRS} (about
4 minutes), saving progress to \texttt{BUILD\_CACHE} after
every 50 sounds, then writes \texttt{bank\_v2.npz} and
\texttt{bank\_v2.csv}. \texttt{NAME\_FIXES} maps the
recording \texttt{tpt-fs3-hmo-mf} (a typo) to
\texttt{tpc-fs3-hmo-mf}.
\texttt{python\ docs/make\_bank\_report.py} then regenerates
the report \texttt{docs/bank\_v2.md} and its figure.

\begin{center}\rule{0.5\linewidth}{0.5pt}\end{center}

\subsection{6. Recipes}

\textbf{A Target from a recording}

\begin{Shaded}
\begin{Highlighting}[]
\ImportTok{from}\NormalTok{ laborch.spectra }\ImportTok{import}\NormalTok{ power\_of\_file}
\NormalTok{target, sus }\OperatorTok{=}\NormalTok{ power\_of\_file(}\StringTok{"my\_target.wav"}\NormalTok{)  }\CommentTok{\# 44.1 kHz; sustain found automatically}
\NormalTok{fit }\OperatorTok{=}\NormalTok{ nnls\_fit(bank.select(instruments}\OperatorTok{=}\NormalTok{[}\StringTok{"flu"}\NormalTok{, }\StringTok{"clb"}\NormalTok{, }\StringTok{"obo"}\NormalTok{]), target)}
\end{Highlighting}
\end{Shaded}

\textbf{Compare the RH2009 methods}

\begin{Shaded}
\begin{Highlighting}[]
\NormalTok{pal }\OperatorTok{=}\NormalTok{ bank.select(instruments}\OperatorTok{=}\NormalTok{[}\StringTok{"flu"}\NormalTok{])}
\NormalTok{t }\OperatorTok{=}\NormalTok{ bank.spectrum(}\StringTok{"vln{-}a4x{-}ord{-}mf"}\NormalTok{)}
\ControlFlowTok{for}\NormalTok{ f }\KeywordTok{in}\NormalTok{ (dot(pal, t, normalize}\OperatorTok{=}\VariableTok{True}\NormalTok{), pinv\_fit(pal, t), nnls\_fit(pal, t)):}
    \BuiltInTok{print}\NormalTok{(f.method, f.ranked(}\DecValTok{3}\NormalTok{, nonzero}\OperatorTok{=}\VariableTok{False}\NormalTok{))}
\end{Highlighting}
\end{Shaded}

\textbf{A synthetic Target} (notebook 2.1, on Bank v1)

\begin{Shaded}
\begin{Highlighting}[]
\ImportTok{from}\NormalTok{ laborch.synth }\ImportTok{import}\NormalTok{ tone}
\NormalTok{bank1 }\OperatorTok{=}\NormalTok{ Bank.load(version}\OperatorTok{=}\DecValTok{1}\NormalTok{)}
\NormalTok{f0 }\OperatorTok{=}\NormalTok{ parse\_name(}\StringTok{"mrb{-}c3{-}ord{-}mf"}\NormalTok{).freq}
\NormalTok{spec }\OperatorTok{=}\NormalTok{ make\_fft(tone(f0, partials}\OperatorTok{=}\NormalTok{[}\DecValTok{1}\NormalTok{, }\DecValTok{4}\NormalTok{, }\DecValTok{10}\NormalTok{], amps}\OperatorTok{=}\NormalTok{[}\DecValTok{1}\NormalTok{, }\FloatTok{0.5}\NormalTok{, }\FloatTok{0.25}\NormalTok{]))}
\BuiltInTok{print}\NormalTok{(nnls\_fit(bank1.select(instruments}\OperatorTok{=}\NormalTok{[}\StringTok{"mrb"}\NormalTok{]), spec).summary(}\DecValTok{3}\NormalTok{))}
\CommentTok{\# mrb{-}c3 0.657, mrb{-}c5 0.298, mrb{-}e6 0.148: a marimba C3 doubled}
\CommentTok{\# on its 4th (C5) and 10th (about E6) partials}
\end{Highlighting}
\end{Shaded}

\textbf{Try a frequency weighting} (relative error per bin)

\begin{Shaded}
\begin{Highlighting}[]
\ImportTok{import}\NormalTok{ numpy }\ImportTok{as}\NormalTok{ np}
\NormalTok{sigma }\OperatorTok{=}\NormalTok{ np.maximum(target, target.}\BuiltInTok{max}\NormalTok{() }\OperatorTok{*} \FloatTok{1e{-}4}\NormalTok{)   }\CommentTok{\# sigma\_i proportional to T\_i}
\NormalTok{fit\_rel }\OperatorTok{=}\NormalTok{ nnls\_fit(palette, target, sigma}\OperatorTok{=}\NormalTok{sigma)}
\end{Highlighting}
\end{Shaded}

\textbf{Select by family, or check the analysis of one
sound}

\begin{Shaded}
\begin{Highlighting}[]
\NormalTok{pal }\OperatorTok{=}\NormalTok{ bank.select(families}\OperatorTok{=}\NormalTok{[}\StringTok{"flutes"}\NormalTok{, }\StringTok{"clarinets"}\NormalTok{], modes}\OperatorTok{=}\StringTok{"ord"}\NormalTok{)}
\NormalTok{row }\OperatorTok{=}\NormalTok{ bank.meta[}\StringTok{"vln{-}a4x{-}ord{-}mf"}\NormalTok{]}
\NormalTok{row[}\StringTok{"sustain\_start\_s"}\NormalTok{], row[}\StringTok{"sustain\_end\_s"}\NormalTok{], row[}\StringTok{"level\_db"}\NormalTok{], row[}\StringTok{"method"}\NormalTok{]}
\end{Highlighting}
\end{Shaded}

\textbf{Plot Target against a Recommendation}

\begin{Shaded}
\begin{Highlighting}[]
\ImportTok{import}\NormalTok{ matplotlib.pyplot }\ImportTok{as}\NormalTok{ plt}
\ImportTok{from}\NormalTok{ laborch.bank }\ImportTok{import}\NormalTok{ frequencies}
\NormalTok{f }\OperatorTok{=}\NormalTok{ frequencies()}\OperatorTok{;}\NormalTok{ r }\OperatorTok{=}\NormalTok{ fit.top(}\DecValTok{4}\NormalTok{, refit}\OperatorTok{=}\VariableTok{True}\NormalTok{)}
\NormalTok{plt.semilogy(f, target, label}\OperatorTok{=}\StringTok{"Target"}\NormalTok{)}\OperatorTok{;}\NormalTok{ plt.semilogy(f, r.model, label}\OperatorTok{=}\StringTok{"Top 4"}\NormalTok{)}
\NormalTok{plt.xlim(}\DecValTok{0}\NormalTok{, }\DecValTok{5000}\NormalTok{)}\OperatorTok{;}\NormalTok{ plt.xlabel(}\StringTok{"Hz"}\NormalTok{)}\OperatorTok{;}\NormalTok{ plt.legend()}
\end{Highlighting}
\end{Shaded}

\textbf{One note per instrument (manual, until Phase 2)} ---
keep the best weight per instrument and refit:

\begin{Shaded}
\begin{Highlighting}[]
\NormalTok{best }\OperatorTok{=}\NormalTok{ \{\}}
\ControlFlowTok{for}\NormalTok{ name, w }\KeywordTok{in}\NormalTok{ fit.ranked():}
\NormalTok{    best.setdefault(bank.meta[name][}\StringTok{"instrument"}\NormalTok{], name)}
\ImportTok{import}\NormalTok{ numpy }\ImportTok{as}\NormalTok{ np}
\ImportTok{from}\NormalTok{ laborch }\ImportTok{import}\NormalTok{ Palette}
\NormalTok{pal1 }\OperatorTok{=}\NormalTok{ Palette(}\BuiltInTok{list}\NormalTok{(best.values()),}
\NormalTok{               np.column\_stack([bank.spectrum(n) }\ControlFlowTok{for}\NormalTok{ n }\KeywordTok{in}\NormalTok{ best.values()]))}
\BuiltInTok{print}\NormalTok{(nnls\_fit(pal1, target).summary())}
\end{Highlighting}
\end{Shaded}

\subsection{7. Data conventions}

\begin{itemize}
\tightlist
\item
  \textbf{Sound names} --- \texttt{inst-pitch-mode-dyn},
  e.g.~\texttt{vln-a4x-ord-mf}. Sharps as \texttt{s}
  (\texttt{cs4}); optional string number before the pitch
  (\texttt{ctb-1b4-nhr-mf}); optional variant after it
  (\texttt{a4o} open string, \texttt{a4x} stopped).
\item
  \textbf{Spectra, Bank v2 (default)} --- 2,048 bins, bin
  \emph{k} at $k \times 10.77$~Hz (bin 0 = DC, last bin $\approx$ 22
  kHz); \textbf{power}, averaged over the detected sustain
  (or the first $\le$ 0.5 s of a struck sound); samples scaled
  to −1..1; each spectrum sums to the mean square of its
  region, so \texttt{10\ log₁₀(sum)} is its level in dBFS.
  Stored as float32.
\item
  \textbf{Spectra, Bank v1} --- same bins; magnitude,
  √(Σ\textbar X\textbar²)/n over the middle 70\% of the
  note, normalized to unit sum.
\item
  \textbf{Levels} --- v2 keeps the \emph{recorded} level,
  which is not the playing level: piano ff is about 4 dB
  quieter than pp in the McGill tracks.
  \texttt{Bank.load(normalize="unit")} removes levels.
\item
  \textbf{Bank coverage} --- 1,515 sounds; nearly all
  \texttt{mf}, with \texttt{pp} and \texttt{ff} only for
  piano; 29 glockenspiel notes listed in the spreadsheets
  were never split out (see \texttt{docs/bank\_v2.md}).
\item
  \textbf{Units of weights} --- v2: a power gain on the
  recording (1 = as recorded). v1: a share of the fitted
  spectrum.
\end{itemize}

\subsection{8. Testing}

\begin{Shaded}
\begin{Highlighting}[]
\BuiltInTok{cd}\NormalTok{ LabOrch}
\ExtensionTok{pytest}\NormalTok{            \# 61 tests, about 10–15 s}
\ExtensionTok{pytest}\NormalTok{ {-}v {-}k ravel}
\end{Highlighting}
\end{Shaded}

\begin{longtable}[]{@{}ll@{}}
\toprule
\begin{minipage}[b]{0.47\columnwidth}\raggedright
Test file\strut
\end{minipage} &
\begin{minipage}[b]{0.47\columnwidth}\raggedright
Checks\strut
\end{minipage}\tabularnewline
\midrule
\endhead
\begin{minipage}[t]{0.47\columnwidth}\raggedright
\texttt{test\_names.py}\strut
\end{minipage} &
\begin{minipage}[t]{0.47\columnwidth}\raggedright
name parsing, MIDI numbers, every Bank name parses (v1 and
v2), v1 unit normalization\strut
\end{minipage}\tabularnewline
\begin{minipage}[t]{0.47\columnwidth}\raggedright
\texttt{test\_metadata.py}\strut
\end{minipage} &
\begin{minipage}[t]{0.47\columnwidth}\raggedright
families, envelope classes, a complete row per sound,
variants on open strings, ranges, v2 regions close to v1,
\texttt{select(families=...)}\strut
\end{minipage}\tabularnewline
\begin{minipage}[t]{0.47\columnwidth}\raggedright
\texttt{test\_power.py}\strut
\end{minipage} &
\begin{minipage}[t]{0.47\columnwidth}\raggedright
Parseval level, powers add, sustain and decay detection on
synthetic notes, fallback, given regions, rebuilding four v2
spectra from their \texttt{.aif} files,
\texttt{normalize="unit"}\strut
\end{minipage}\tabularnewline
\begin{minipage}[t]{0.47\columnwidth}\raggedright
\texttt{test\_solve.py}\strut
\end{minipage} &
\begin{minipage}[t]{0.47\columnwidth}\raggedright
NNLS recovers an exact synthetic mixture, $\sigma$ weighting,
\texttt{top}/refit, dB gains\strut
\end{minipage}\tabularnewline
\begin{minipage}[t]{0.47\columnwidth}\raggedright
\texttt{test\_spectra.py}\strut
\end{minipage} &
\begin{minipage}[t]{0.47\columnwidth}\raggedright
\texttt{make\_fft} rebuilds four stored v1 spectra from
their \texttt{.aif} files; pure-tone peak\strut
\end{minipage}\tabularnewline
\begin{minipage}[t]{0.47\columnwidth}\raggedright
\texttt{test\_regression\_notebook\_1\_5.py}\strut
\end{minipage} &
\begin{minipage}[t]{0.47\columnwidth}\raggedright
violin A4 from flutes matches notebook 1.5 to three
decimals; pseudo-inverse goes negative, NNLS does not\strut
\end{minipage}\tabularnewline
\begin{minipage}[t]{0.47\columnwidth}\raggedright
\texttt{test\_synth.py}\strut
\end{minipage} &
\begin{minipage}[t]{0.47\columnwidth}\raggedright
\texttt{tone} partials, amplitudes and decay; the synthetic
C3 decomposes into marimba C3 + C5 + E6; spectrum-to-audio
round trip\strut
\end{minipage}\tabularnewline
\begin{minipage}[t]{0.47\columnwidth}\raggedright
\texttt{test\_regression\_ravel.py}\strut
\end{minipage} &
\begin{minipage}[t]{0.47\columnwidth}\raggedright
RH2009 \emph{Daphnis et Chloé} example, on v1 and v2: the
woodwind chord is closer to the cup mute than the straight
mute\strut
\end{minipage}\tabularnewline
\bottomrule
\end{longtable}

Audio tests are skipped automatically if pydub or ffmpeg is
missing; data tests are skipped if the Bank folders are
absent.

\subsection{9. Known limitations and roadmap}

\begin{longtable}[]{@{}ll@{}}
\toprule
\begin{minipage}[b]{0.47\columnwidth}\raggedright
Limitation\strut
\end{minipage} &
\begin{minipage}[b]{0.47\columnwidth}\raggedright
Planned in\strut
\end{minipage}\tabularnewline
\midrule
\endhead
\begin{minipage}[t]{0.47\columnwidth}\raggedright
Recorded levels are not playing levels (piano ff $\approx$ 4 dB
below pp); no loudness model\strut
\end{minipage} &
\begin{minipage}[t]{0.47\columnwidth}\raggedright
Phase 2 (or new recordings)\strut
\end{minipage}\tabularnewline
\begin{minipage}[t]{0.47\columnwidth}\raggedright
Every frequency bin counts equally ($\sigma = 1$); with power
spectra the loudest partials dominate\strut
\end{minipage} &
\begin{minipage}[t]{0.47\columnwidth}\raggedright
Phase 2 (choice of $\sigma$)\strut
\end{minipage}\tabularnewline
\begin{minipage}[t]{0.47\columnwidth}\raggedright
No playability: several notes may be assigned to one
player\strut
\end{minipage} &
\begin{minipage}[t]{0.47\columnwidth}\raggedright
Phase 2 (one-note-per-player search)\strut
\end{minipage}\tabularnewline
\begin{minipage}[t]{0.47\columnwidth}\raggedright
Dynamics are mf only (except piano); 29 glockenspiel notes
missing; bass sax has 8 notes\strut
\end{minipage} &
\begin{minipage}[t]{0.47\columnwidth}\raggedright
Phase 1 follow-up: re-split the glockenspiel track; source
pp/ff recordings\strut
\end{minipage}\tabularnewline
\begin{minipage}[t]{0.47\columnwidth}\raggedright
No web interface, score export, morphing or creative
mode\strut
\end{minipage} &
\begin{minipage}[t]{0.47\columnwidth}\raggedright
Phases 3--5\strut
\end{minipage}\tabularnewline
\bottomrule
\end{longtable}

The full plan is in the project's Planning documents (not in
the public repo).

\end{document}
